% notation.tex --- Core notation for
% "Odd-Rank Loci and Box-Induced Symmetries of Centered Sudoku Operators".
%
% This snippet is self-contained and is \input by every section that
% needs the symbols. It is consistent with the sibling paper
% "Determinant Divisibility of Centered Latin Squares".
%
% -----------------------------------------------------------------------
% Macros
% -----------------------------------------------------------------------
\providecommand{\one}{\mathbf{1}}
\providecommand{\Z}{\mathbb{Z}}
\providecommand{\Q}{\mathbb{Q}}
\providecommand{\R}{\mathbb{R}}
\providecommand{\rank}{\operatorname{rank}}
\providecommand{\Aut}{\operatorname{Aut}}
\providecommand{\Stab}{\operatorname{Stab}}
\providecommand{\SNF}{\operatorname{SNF}}

% -----------------------------------------------------------------------
% Setup
% -----------------------------------------------------------------------
% Let $L$ be a Latin square of order $n$. In the main Sudoku case
% $n = 9$ and $L$ is a Sudoku grid (i.e.\ a Latin square also satisfying
% the $3 \times 3$ box constraint). A \emph{symbol relabeling} is a
% permutation $\gamma \in S_n$, applied entrywise to $L$ to produce
% $\gamma(L)$.

% Centering scalar and centered operator:
% \[
%   m_n \;:=\; \frac{n+1}{2},
%   \qquad
%   E_{\gamma,L} \;:=\; \gamma(L) - m_n\, J_n,
% \]
% where $J_n$ is the $n \times n$ all-ones matrix. When $L$ is fixed
% we write $E_\gamma$. For $n = 9$, $m_n = 5$, hence
% $E_\gamma = \gamma(L) - 5J$.

% -----------------------------------------------------------------------
% Standard hyperplane and root lattice
% -----------------------------------------------------------------------
% Standard hyperplane (\emph{vector space}, not a lattice):
% \[
%   V_{\mathrm{std},n} \;:=\; \{\, x \in \Q^n \;:\; \one^\top x = 0 \,\}.
% \]
% Root lattice (the integer points of $V_{\mathrm{std},n}$):
% \[
%   A_{n-1} \;:=\; V_{\mathrm{std},n} \cap \Z^n.
% \]
% In particular, for $n = 9$, $A_8 = V_{\mathrm{std},9} \cap \Z^9$.

% -----------------------------------------------------------------------
% Rank, restriction, and Gram-projected representative
% -----------------------------------------------------------------------
% Since $E_{\gamma,L}\,\one = \mathbf{0}$ and $\one^\top E_{\gamma,L} = \mathbf{0}$,
% the operator $E_{\gamma,L}$ restricts to an endomorphism of
% $V_{\mathrm{std},n}$. Consequently
% \[
%   \rank(E_{\gamma,L})
%   \;=\; \rank\bigl(E_{\gamma,L}\bigr|_{V_{\mathrm{std},n}}\bigr).
% \]
% For SNF (Smith normal form) statements we use the Gram-projected
% representative in the standard basis $P_n = (e_i - e_n)_{i=1}^{n-1}$,
% \[
%   \widehat{E}_{\gamma,L} \;:=\; P_n^{\top}\, E_{\gamma,L}\, P_n,
% \]
% as in the sibling paper. \textbf{Convention.} Rank statements may be
% read on $E_{\gamma,L}$, on its restriction to $V_{\mathrm{std},n}$, or
% on $\widehat{E}_{\gamma,L}$; SNF statements default to
% $\widehat{E}_{\gamma,L}$ unless explicitly stated otherwise.

% -----------------------------------------------------------------------
% Odd-rank locus
% -----------------------------------------------------------------------
% General order:
% \[
%   \Gamma_L^{\mathrm{odd}}
%   \;:=\; \{\, \gamma \in S_n \;:\; \rank(E_{\gamma,L}) \text{ is odd} \,\}.
% \]
% Sudoku case ($n = 9$, where the maximal rank on $V_{\mathrm{std},9}$
% is $8$):
% \[
%   \Gamma_L^{*}
%   \;:=\; \{\, \gamma \in S_9 \;:\; \rank(E_{\gamma,L}) = 7 \,\}
%   \;\subseteq\; \Gamma_L^{\mathrm{odd}}.
% \]

% -----------------------------------------------------------------------
% Kernel certificates
% -----------------------------------------------------------------------
% A \emph{kernel certificate} for $(L, \gamma)$ is a vector
% $u \in A_{n-1}$ with $E_{\gamma,L}\,u = \mathbf{0}$ (right certificate)
% or $u^\top E_{\gamma,L} = \mathbf{0}$ (left certificate). For the
% laboratory base $L = L_{M_{19}}$ at $n = 9$, $u_L^* \in A_8$ denotes
% the dominant shared certificate in the sense of Section~\ref{sec:M19}.

% -----------------------------------------------------------------------
% Affine coset symmetry, autotopisms, template closure
% -----------------------------------------------------------------------
% The groups $K_L, K_R$ are the \emph{left and right affine coset-symmetry
% groups of the bilateral coset structure of $\Gamma^*_{M_{19}}$}
% (cf.\ Theorems in Section~\ref{sec:M19}); they are \textbf{not}
% defined as literal stabilizers of $u_L^*$, and they are
% \textbf{not} Latin-square autotopism groups of $L$.
%
% $\Aut(L)$ denotes the (Latin-square) autotopism group of $L$;
% paratopisms are excluded unless explicitly stated.
%
% $G^*$ denotes the corrected closed group acting on certificate
% templates after the S9w-audit, with side labels (left/right) tracked
% separately.
